\documentclass[12pt,a4paper,openany]{report}
\usepackage[utf8]{inputenc}
\usepackage[T1]{fontenc}
\usepackage{times}              % Times New Roman (equivalente); com XeLaTeX pode usar fontspec
\usepackage[portuguese]{babel}
\usepackage[a4paper,left=2.5cm,top=2.5cm,right=3cm,bottom=3cm]{geometry}
\usepackage{setspace}
\usepackage{titlesec}
\usepackage{tocloft}
\usepackage{caption}
\usepackage{booktabs,tabularx,longtable,array}
\usepackage{enumitem}
\usepackage{graphicx,adjustbox}
\usepackage{tikz}
\usetikzlibrary{positioning,arrows.meta,shapes.geometric,shapes.multipart,fit,calc,backgrounds}
\usepackage{microtype}
\usepackage[hidelinks]{hyperref}

%% ---------- Formatação exigida ----------
\onehalfspacing                                 % espaçamento 1,5
\setlength{\parindent}{1.25cm}
\setlength{\parskip}{0pt}
\titleformat{\chapter}[hang]{\fontsize{14}{17}\selectfont\bfseries}{\MakeUppercase{Capítulo \Roman{chapter}} --}{.5em}{\MakeUppercase}
\titleformat{name=\chapter,numberless}[hang]{\fontsize{14}{17}\selectfont\bfseries}{}{0pt}{\MakeUppercase}
\titlespacing*{\chapter}{0pt}{0pt}{16pt}
\titleformat{\section}{\fontsize{12}{14}\selectfont\bfseries}{\thesection.}{.5em}{}
\titleformat{\subsection}{\fontsize{12}{14}\selectfont\bfseries\itshape}{\thesubsection.}{.5em}{}
\titleformat{\subsubsection}{\fontsize{12}{14}\selectfont\itshape}{\thesubsubsection.}{.5em}{}
\titlespacing*{\section}{0pt}{8pt}{6pt}
\titlespacing*{\subsection}{0pt}{8pt}{4pt}
\titlespacing*{\subsubsection}{0pt}{8pt}{4pt}
\captionsetup{font=small,labelfont=bf,labelsep=period}
\captionsetup[table]{position=above}
\captionsetup[figure]{position=below}
\renewcommand{\cftchapfont}{\bfseries}
\setlength{\cftbeforechapskip}{6pt}
\renewcommand{\arraystretch}{1.15}
\setlist{nosep}

%% ---------- Macros ----------
\newcommand{\ucspec}[8]{% id, nome, actores, descrição, pré, pós, fluxo, alternativos
\begin{table}[!htbp]
\caption{Especificação do caso de uso #1 --- #2}
\label{tab:#1}
\begingroup\singlespacing\small
\begin{tabularx}{\textwidth}{@{}>{\bfseries}p{3cm}X@{}}
\toprule
Actores & #3 \\
Descrição & #4 \\
Pré-condições & #5 \\
Pós-condições & #6 \\
Fluxo principal & #7 \\
Fluxos alternativos & #8 \\
\bottomrule
\end{tabularx}
\endgroup
\end{table}}

\newcommand{\autorA}{kleyton Aquimo Manuel}
\newcommand{\autorB}{Silvano Filipe Jose}
\newcommand{\autorC}{Valther Antonio Canze\\Yuren Fernando Franco}


\begin{document}
\pagenumbering{roman}

%% ---------- CAPA ----------
\begin{titlepage}
\centering\singlespacing
{\fontsize{14}{17}\selectfont\bfseries Universidade zambeze\par}
{\fontsize{12}{14}\selectfont Faculdade de ciencias e tecnologias\par}
\vspace{3.5cm}
{\fontsize{14}{17}\selectfont\bfseries SISTEMA INTEGRADO DE INCLUSÃO FINANCEIRA MUNICIPAL (SIIFM):\\[4pt] MODELAGEM E CONCEPÇÃO\par}
\vspace{1.2cm}
{\fontsize{12}{14}\selectfont Relatório da disciplina de Sistemas de informacao\par}
\vfill

% Estudantes à Esquerda
\begin{flushleft}\fontsize{12}{14}\selectfont
\textbf{Estudantes:}\\ \autorA\\ \autorB\\ \autorC
\end{flushleft}

\vspace{1cm} % Espaço entre estudantes e docente

% Docente à Direita
\begin{flushright}\fontsize{12}{14}\selectfont
\textbf{Docente:} eng Florencio Capachiua
\end{flushright}

\vfill
{\fontsize{12}{14}\selectfont Beira, Outubro de 2026\par}
\end{titlepage}

%% ---------- RESUMO ----------
\chapter*{Resumo}
\addcontentsline{toc}{chapter}{Resumo}
Os municípios moçambicanos dependem de taxas e licenças cobradas, em grande parte, de forma manual, em numerário e com recibos em papel, o que dificulta o controlo da receita, a rastreabilidade das operações e o acesso do cidadão aos serviços. Em paralelo, as carteiras móveis (M-Pesa e e-Mola) tornaram-se um canal financeiro amplamente utilizado. Este trabalho apresenta a modelagem e a concepção do Sistema Integrado de Inclusão Financeira Municipal (SIIFM), uma plataforma que permite registar contribuintes, pagar taxas municipais por carteira móvel, POS ou numerário, emitir recibos digitais e fornecer relatórios de arrecadação aos gestores municipais.

A abordagem seguiu o ciclo de vida de desenvolvimento de software (SDLC) com práticas ágeis (Scrum) e a notação UML. Os resultados incluem o levantamento de requisitos funcionais e não funcionais, oito casos de uso, o diagrama de classes, os diagramas de sequência e de actividades, o modelo entidade-relacionamento, a arquitectura em camadas, o plano de testes, o plano de sprints, a análise de riscos e e a implementação do sistema (API em Node.js, base de dados SQL e interface web com portais separados para o contribuinte, o agente de balcão, o gestor e o administrador), verificada por testes automáticos. Conclui-se que o modelo proposto é coerente e viável, ficando a validação empírica com utilizadores reais e a integração com as APIs reais das operadoras como trabalho futuro.

\vspace{6pt}\noindent\textbf{Palavras-chave:} inclusão financeira; carteiras móveis; taxas municipais; modelagem de sistemas; UML.

%% ---------- LISTAS ----------
\cleardoublepage\clearpage
\renewcommand{\listfigurename}{Lista de Figuras}
\renewcommand{\listtablename}{Lista de Tabelas}
\renewcommand{\contentsname}{Sumário}
\begingroup
\titleformat{\chapter}[hang]{\fontsize{14}{17}\selectfont\bfseries}{}{0pt}{\MakeUppercase}
\listoffigures
\addcontentsline{toc}{chapter}{Lista de Figuras}
\clearpage
\listoftables
\addcontentsline{toc}{chapter}{Lista de Tabelas}
\clearpage
\tableofcontents
\endgroup

\clearpage
\pagenumbering{arabic}

%% =====================================================
\chapter{Introdução}
Os governos locais são a ponta da administração pública mais próxima do cidadão e do pequeno comerciante. É no município que se pagam as taxas de mercado, as licenças de comércio e a taxa de saneamento, e é dessa receita própria que depende parte da capacidade de o município prestar serviços. Quando a cobrança é manual, a receita fica exposta a perdas, atrasos e erros de registo, e o contribuinte perde tempo em deslocações e filas.

Ao mesmo tempo, o dinheiro móvel alterou a forma como as pessoas guardam, enviam e pagam dinheiro em Moçambique e noutros países africanos \cite{jack2011,findex2021}. Este contexto cria uma oportunidade: utilizar as carteiras móveis já existentes como canal para a arrecadação municipal, promovendo simultaneamente a inclusão financeira \cite{enif}. O presente relatório documenta a modelagem e a concepção do Sistema Integrado de Inclusão Financeira Municipal (SIIFM) com esse propósito.

\section{Problematização}
A arrecadação de taxas municipais enfrenta, de forma recorrente, os seguintes problemas:
\begin{itemize}
\item \textbf{Processos manuais e dispersos:} registos em papel ou em folhas de cálculo isoladas, sem perfil único do contribuinte e com risco de duplicação de dados;
\item \textbf{Pagamentos em numerário sem rastreio automático:} dificultam a conciliação, aumentam o risco de desvio e tornam a auditoria lenta;
\item \textbf{Barreiras de acesso:} o contribuinte tem de se deslocar aos serviços municipais, em horário limitado, para pagar e levantar o comprovativo;
\item \textbf{Falta de informação de gestão:} os relatórios de arrecadação são produzidos tardiamente, o que limita o planeamento e a decisão;
\item \textbf{Baixa inclusão financeira digital dos pequenos comerciantes:} apesar da adopção do dinheiro móvel, as taxas municipais raramente podem ser pagas por esse canal.
\end{itemize}

Perante este quadro, formula-se a seguinte \textbf{pergunta de investigação}: \emph{como conceber um sistema integrado que permita registar contribuintes, cobrar taxas municipais através de carteiras móveis e de outros meios, emitir recibos digitais e fornecer relatórios fiáveis ao município, com segurança, rastreabilidade e acessibilidade?}

\section{Hipóteses}
\begin{description}[leftmargin=1.5cm,style=nextline]
\item[H1.] A disponibilização do pagamento de taxas por carteira móvel, POS e numerário num único sistema reduz o esforço do contribuinte e as falhas de registo da receita.
\item[H2.] A centralização do perfil do contribuinte, do histórico de pagamentos e do registo de auditoria melhora a rastreabilidade e a transparência da arrecadação.
\item[H3.] Uma arquitectura web/PWA com canais alternativos (SMS e USSD) permite o acesso de utilizadores com conectividade limitada.
\item[H4.] Relatórios e \emph{dashboards} em tempo real melhoram o apoio à decisão dos gestores municipais.
\end{description}

\section{Objetivos}
\subsection{Objetivo Geral}
Modelar e conceber o Sistema Integrado de Inclusão Financeira Municipal (SIIFM) para a gestão digital de contribuintes, o pagamento de taxas municipais e a produção de relatórios de arrecadação.

\subsection{Objetivos específicos}
\begin{enumerate}
\item Identificar os actores, os requisitos funcionais e os requisitos não funcionais do sistema;
\item Modelar o comportamento do sistema através de casos de uso, diagramas de sequência e de actividades;
\item Modelar a estrutura do sistema através dos diagramas de classes e do modelo entidade-relacionamento;
\item Definir a arquitectura técnica, as medidas de segurança e a conformidade com a protecção de dados;
\item Planear o desenvolvimento (Scrum), os testes e a gestão de riscos;
\item Implementar um sistema funcional (API, base de dados e interface) com portais para o contribuinte, o agente de balcão, o gestor e o administrador, e verificá-lo com testes automáticos.
\end{enumerate}

\section{Organização do Trabalho}
O relatório está organizado em quatro capítulos. O \textbf{Capítulo I} apresenta o contexto, a problematização, as hipóteses e os objectivos. O \textbf{Capítulo II} reúne o marco teórico sobre inclusão financeira, engenharia de software, Scrum, UML, arquitectura web, segurança e testes. O \textbf{Capítulo III} apresenta os resultados e a modelagem do SIIFM: requisitos, casos de uso, diagramas, arquitectura, plano de testes, plano de sprints, riscos e o sistema implementado. O \textbf{Capítulo IV} apresenta as conclusões, as limitações e as recomendações. Seguem-se a bibliografia e os apêndices.

%% =====================================================
\chapter{Marco Teórico}

\section{Inclusão financeira e dinheiro móvel}
A inclusão financeira consiste no acesso e na utilização, por indivíduos e empresas, de serviços financeiros úteis e acessíveis, como pagamentos, poupança e crédito \cite{findex2021}. Em países com baixa penetração bancária, o dinheiro móvel permitiu que milhões de pessoas efectuassem pagamentos digitais sem conta bancária; o caso do M-PESA, no Quénia, é a referência mais estudada \cite{jack2011}. Em Moçambique, a Estratégia Nacional de Inclusão Financeira reconhece os canais digitais como instrumento de alargamento do acesso \cite{enif}. As carteiras M-Pesa e e-Mola são, neste trabalho, os principais meios de pagamento digital considerados.

\section{Gestão tributária municipal e governo electrónico}
O governo electrónico refere-se ao uso de tecnologias de informação para melhorar a prestação de serviços públicos, a eficiência interna e a transparência. No domínio municipal, a digitalização da cobrança implica (i) registo unificado de contribuintes, (ii) catálogo de taxas com valor e periodicidade, (iii) emissão de comprovativo e (iv) informação de gestão. A informatização não elimina a necessidade de canais presenciais (numerário e POS), que continuam a ser indispensáveis para quem não tem carteira móvel; por isso o SIIFM é concebido como sistema \emph{multicanal}.

\section{Engenharia de software e ciclo de vida (SDLC)}
A engenharia de software é a aplicação de princípios sistemáticos ao desenvolvimento, à operação e à manutenção de software \cite{sommerville2016,pressman2020}. O ciclo de vida de desenvolvimento (SDLC) organiza o trabalho nas fases de requisitos, \emph{design}, implementação, testes, \emph{deploy} e manutenção. A fase de requisitos distingue os \textbf{requisitos funcionais}, que descrevem o que o sistema faz, dos \textbf{requisitos não funcionais}, que descrevem propriedades como desempenho, segurança e disponibilidade \cite{sommerville2016}.

\section{Metodologia ágil Scrum}
O Scrum é uma estrutura leve para o desenvolvimento incremental de produtos complexos \cite{scrum2020}. Organiza o trabalho em \emph{sprints} de duração fixa, nos quais uma equipa multifuncional entrega um incremento utilizável. Os artefactos centrais são o \emph{Product Backlog} (lista priorizada por valor de negócio), o \emph{Sprint Backlog} e o incremento; os eventos são o planeamento da \emph{sprint}, o \emph{Daily Scrum}, a revisão e a retrospectiva. Este modelo é adequado ao SIIFM por permitir demonstrações regulares aos \emph{stakeholders} municipais e a incorporação do seu \emph{feedback}.

\section{Linguagem de modelagem unificada (UML)}
A UML é a linguagem padrão para especificar, visualizar e documentar sistemas de software \cite{booch2005,fowler2004}. Neste trabalho utilizam-se os seguintes diagramas:
\begin{itemize}
\item \textbf{Casos de uso:} funcionalidades do sistema do ponto de vista dos actores;
\item \textbf{Classes:} estrutura estática, atributos, operações e associações;
\item \textbf{Sequência:} troca temporal de mensagens entre os participantes de um cenário;
\item \textbf{Actividades:} fluxo de trabalho com decisões e ramificações.
\end{itemize}
A análise orientada a objectos e a derivação de casos de uso seguem a abordagem descrita por Larman \cite{larman2004}.

\section{Arquitectura web: SPA, PWA e API REST}
Uma \emph{Single Page Application} (SPA) carrega uma única página e actualiza a interface dinamicamente, comunicando com o servidor através de uma API. Uma \emph{Progressive Web App} (PWA) acrescenta capacidades como instalação no ecrã inicial e funcionamento parcialmente \emph{offline}, útil em contextos de conectividade intermitente. O estilo arquitectural REST \cite{fielding2000} baseia-se em recursos identificados por URL, operações HTTP padronizadas e ausência de estado no servidor, o que facilita a escalabilidade horizontal e a integração com sistemas externos como as operadoras móveis.

\section{Bases de dados relacionais}
O modelo relacional organiza os dados em tabelas ligadas por chaves primárias e estrangeiras, garantindo integridade referencial e propriedades transaccionais ACID \cite{elmasri2016}. Estas propriedades são essenciais para pagamentos, em que a gravação do pagamento, do recibo e do registo de auditoria deve ocorrer de forma atómica. O modelo entidade-relacionamento (ER) é usado na concepção lógica, e a normalização reduz a redundância.

\section{Segurança e protecção de dados}
Os sistemas que tratam dados pessoais e financeiros devem aplicar encriptação em trânsito (TLS) e em repouso, controlo de acessos por perfil (RBAC), armazenamento de credenciais com \emph{hash} e \emph{salt}, registo de auditoria e protecção contra as vulnerabilidades mais frequentes, tais como as descritas no OWASP Top 10 \cite{owasp2021}. No plano jurídico, o enunciado do projecto exige conformidade com o RGPD \cite{rgpd2016} e com a legislação moçambicana aplicável, incluindo a Lei das Transacções Electrónicas \cite{lei2017}, cuja aplicação concreta a dados pessoais deve ser confirmada pelo jurista da instituição antes da operação real.

\section{Testes de software}
O teste de software visa detectar defeitos e verificar se o sistema cumpre os requisitos \cite{myers2011}. Distinguem-se testes \textbf{unitários} (módulos isolados), de \textbf{integração} (interacção entre módulos), \textbf{ponta a ponta} (E2E, fluxo completo), de \textbf{usabilidade} \cite{nielsen1993} e de \textbf{acessibilidade} (WCAG \cite{wcag2018}), de \textbf{segurança} (incluindo testes de intrusão) e de \textbf{desempenho} (carga e latência).

\section{Síntese do capítulo}
O marco teórico mostra que o problema pode ser tratado combinando (i) o dinheiro móvel como canal de inclusão, (ii) um processo de engenharia disciplinado (SDLC com Scrum), (iii) a modelagem UML para comunicar o desenho e (iv) uma arquitectura web em camadas com segurança e testes desde o início. Estes elementos orientam o Capítulo III.
%% =====================================================
\chapter{Apresentação dos Resultados e Modelagem}

\section{Abordagem metodológica}
A modelagem seguiu o SDLC, com as fases de requisitos e \emph{design} documentadas neste relatório e as restantes planeadas por \emph{sprints} (Secção~\ref{sec:scrum}). Partiu-se da especificação inicial do SIIFM (quatro casos de uso, arquitectura técnica, sete requisitos funcionais, seis não funcionais, metodologia e tipos de teste), que foi analisada, completada e formalizada em UML. O sistema implementado (Secção~\ref{sec:proto}) foi construído e testado para validar os principais fluxos. Os requisitos e os casos de uso marcados como \emph{derivados} foram acrescentados pela equipa por serem necessários para tornar o sistema coerente (por exemplo, autenticação e consulta do histórico).

\section{Actores e partes interessadas}
\begin{table}[!htbp]
\caption{Actores do sistema}\label{tab:actores}
\begingroup\singlespacing\small
\begin{tabularx}{\textwidth}{@{}>{\bfseries}p{3.6cm}X@{}}
\toprule
Actor & Papel \\
\midrule
Cidadão & Pessoa singular que se regista, paga taxas (por exemplo, a taxa de saneamento) e consulta recibos. \\
Comerciante & Contribuinte com actividade comercial; paga licenças e taxas de mercado. Partilha as funcionalidades do cidadão. \\
Administrador municipal & Gere taxas e categorias, consulta contribuintes, relatórios e auditoria. \\
Gestor municipal & Consulta relatórios e \emph{dashboards} para apoio à decisão. \\
Operadora móvel (M-Pesa, e-Mola) & Sistema externo que autoriza o débito da carteira e confirma a transacção. \\
Serviço de notificações & Sistema externo de SMS e e-mail para confirmações e recibos. \\
Agente municipal (POS/numerário) & Funcionário que regista pagamentos presenciais em POS ou numerário. \\
\bottomrule
\end{tabularx}
\endgroup
\end{table}

\section{Requisitos funcionais}
Os requisitos RF01 a RF07 provêm da especificação inicial; RF08 e RF09 são derivados.

\begin{table}[!htbp]
\caption{Requisitos funcionais}\label{tab:rf}
\begingroup\singlespacing\small
\begin{tabularx}{\textwidth}{@{}>{\bfseries}p{1.3cm}Xp{2.2cm}@{}}
\toprule
ID & Descrição & Prioridade \\
\midrule
RF01 & Registo e autenticação de utilizadores (cidadão, comerciante, administrador). & Alta \\
RF02 & Pagamento de taxas por carteiras móveis e POS (e numerário). & Alta \\
RF03 & Histórico completo de pagamentos e emissão de recibos digitais. & Alta \\
RF04 & Gestão de taxas, categorias e periodicidade pelo administrador. & Alta \\
RF05 & Relatórios e \emph{dashboards} em tempo real; exportação em PDF/Excel. & Média \\
RF06 & Notificações e alertas por SMS e e-mail. & Média \\
RF07 & Auditoria de todas as transacções e alterações. & Alta \\
RF08 & Detecção de duplicação de contribuintes no registo (derivado). & Média \\
RF09 & Acesso por múltiplos canais: web, app/PWA e USSD (derivado). & Média \\
\bottomrule
\end{tabularx}
\endgroup
\end{table}

\section{Requisitos não funcionais}
\begin{table}[!htbp]
\caption{Requisitos não funcionais e forma de verificação}\label{tab:rnf}
\begingroup\singlespacing\small
\begin{tabularx}{\textwidth}{@{}>{\bfseries}p{1.5cm}p{2.6cm}Xp{3cm}@{}}
\toprule
ID & Categoria & Requisito & Verificação \\
\midrule
RNF01 & Disponibilidade & Disponibilidade mínima de 99\%. & Monitorização \\
RNF02 & Desempenho & Latência inferior a 2~s nas acções críticas (pagamento, registo). & Teste de carga \\
RNF03 & Segurança & TLS~1.3 e encriptação de dados sensíveis. & Teste de segurança \\
RNF04 & Escalabilidade & Suportar até 50\,000 utilizadores. & Teste de carga \\
RNF05 & Usabilidade & Interface acessível e intuitiva (WCAG básico). & Teste com utilizadores \\
RNF06 & Conformidade & RGPD e legislação moçambicana de protecção de dados. & Revisão jurídica \\
RNF07 & Recuperação & Cópias de segurança automáticas e \emph{logs} auditáveis (derivado). & Teste de restauro \\
\bottomrule
\end{tabularx}
\endgroup
\end{table}

\section{Modelagem de casos de uso}
O diagrama da Figura~\ref{fig:uc} apresenta os oito casos de uso do SIIFM e os seus actores. UC01 a UC04 provêm da especificação inicial; UC05 a UC08 são derivados.

\begin{figure}[!htbp]
\centering
\begin{adjustbox}{max width=\textwidth}
\begin{tikzpicture}[font=\small,
  uc/.style={ellipse,draw,minimum width=4.6cm,minimum height=.9cm,align=center,inner sep=1pt},
  act/.style={draw,rounded corners,minimum width=2.4cm,minimum height=.9cm,align=center,fill=gray!10}]
\draw[rounded corners] (3.2,1) rectangle (11.2,-11.8);
\node at (7.2,.6) {\textbf{Sistema SIIFM}};
\node[uc] (u1) at (7.2,-.5) {UC01 Registar contribuinte};
\node[uc] (u5) at (7.2,-2) {UC05 Autenticar-se};
\node[uc] (u2) at (7.2,-3.5) {UC02 Pagar taxa municipal};
\node[uc] (u6) at (7.2,-5) {UC06 Consultar histórico e recibos};
\node[uc] (u3) at (7.2,-6.5) {UC03 Gerir taxas e categorias};
\node[uc] (u4) at (7.2,-8) {UC04 Gerar relatórios e dashboards};
\node[uc] (u7) at (7.2,-9.5) {UC07 Consultar auditoria};
\node[uc] (u8) at (7.2,-11) {UC08 Registar pagamento presencial};
\node[act] (a1) at (0,-3) {Cidadão /\\Comerciante};
\node[act] (a2) at (14.4,-7.2) {Administrador /\\Gestor};
\node[act] (a3) at (14.4,-3.5) {Operadora móvel\\(M-Pesa, e-Mola)};
\node[act] (a4) at (14.4,-0.5) {Serviço de\\notificações};
\node[act] (a5) at (0,-10) {Agente\\municipal};
\draw (a1.east) -- (u1.west); \draw (a1.east) -- (u5.west); \draw (a1.east) -- (u2.west); \draw (a1.east) -- (u6.west);
\draw (a2.west) -- (u5.east); \draw (a2.west) -- (u3.east); \draw (a2.west) -- (u4.east); \draw (a2.west) -- (u7.east);
\draw (a3.west) -- (u2.east);
\draw (a4.west) -- (u1.east); \draw (a4.west) -- (u2.east);
\draw (a5.east) -- (u8.west); \draw (a5.east) -- (u5.west);
\end{tikzpicture}
\end{adjustbox}
\caption{Diagrama de casos de uso do SIIFM. UC07 e UC08 estão resumidos na Tabela~\ref{tab:UC08}.}\label{fig:uc}
\end{figure}

\ucspec{UC01}{Registo de contribuinte}
{Cidadão ou comerciante (principal); serviço de notificações (secundário).}
{Permite registar-se no sistema com dados pessoais e contacto, criando um perfil único.}
{O utilizador ainda não possui perfil e tem acesso à aplicação (web, app ou USSD).}
{Perfil criado com PIN em \emph{hash}; confirmação enviada por SMS ou e-mail; evento gravado na auditoria.}
{\begin{enumerate}[leftmargin=1.2em]
\item O utilizador acede à aplicação;
\item Preenche nome, contacto, morada e tipo (cidadão ou comerciante) e define um PIN;
\item O sistema valida os dados e verifica duplicação;
\item O sistema cria o perfil;
\item O sistema envia a confirmação por SMS ou e-mail.
\end{enumerate}}
{A1: contacto já registado: o sistema alerta duplicação e não cria perfil. A2: dados inválidos (telemóvel fora do formato, PIN curto): o sistema indica o campo a corrigir.}

\ucspec{UC02}{Pagamento de taxas municipais}
{Cidadão ou comerciante (principal); operadora móvel e serviço de notificações (secundários).}
{Permite pagar taxas por carteira móvel, POS ou numerário, com emissão de recibo digital.}
{O utilizador está autenticado e existe pelo menos uma taxa activa.}
{Pagamento registado com estado \emph{confirmado}; recibo emitido; histórico actualizado; auditoria gravada.}
{\begin{enumerate}[leftmargin=1.2em]
\item O utilizador selecciona a taxa ou licença e o número de períodos;
\item O sistema calcula o valor e gera a referência de pagamento;
\item O utilizador escolhe o método (carteira móvel, POS ou numerário) e paga;
\item O sistema valida a confirmação e grava o pagamento;
\item O sistema actualiza o histórico e envia o recibo digital.
\end{enumerate}}
{A1: a operadora recusa ou expira o pedido: o pagamento fica \emph{falhado} e o utilizador pode repetir. A2: perda de ligação após o débito: o sistema reconcilia o estado pelo \emph{callback} da operadora antes de permitir novo débito.}

\ucspec{UC03}{Gestão de taxas pelo administrador}
{Administrador municipal.}
{Permite criar, alterar e eliminar taxas, bem como gerir categorias.}
{O administrador está autenticado com perfil de administrador.}
{Taxa disponível (ou removida) para pagamento; alteração gravada na auditoria.}
{\begin{enumerate}[leftmargin=1.2em]
\item O administrador autentica-se;
\item Cria uma nova taxa com nome, valor, periodicidade e categoria;
\item O sistema valida e disponibiliza a taxa para pagamento.
\end{enumerate}}
{A1: valor inválido (zero ou negativo): o sistema rejeita. A2: eliminação de taxa com pagamentos associados: o sistema desactiva a taxa em vez de a apagar, preservando o histórico.}

\ucspec{UC04}{Relatórios e \emph{dashboards}}
{Administrador ou gestor municipal.}
{Permite gerar relatórios diários, semanais e mensais de arrecadação.}
{O utilizador está autenticado com perfil de administrador ou gestor.}
{Relatório apresentado e exportado; consulta registada na auditoria.}
{\begin{enumerate}[leftmargin=1.2em]
\item O utilizador seleciona o período e o tipo de relatório;
\item O sistema gera os dados (por taxa, método e período);
\item O sistema apresenta e permite exportar em PDF ou Excel.
\end{enumerate}}
{A1: período sem pagamentos: o sistema mostra o relatório vazio com mensagem explicativa.}

\ucspec{UC05}{Autenticação (derivado)}
{Cidadão, comerciante, administrador ou gestor.}
{Permite entrar no sistema com contacto e PIN/palavra-passe e encaminha para o portal correspondente ao perfil.}
{O utilizador possui perfil activo.}
{Sessão iniciada com o perfil correcto; entrada registada na auditoria.}
{\begin{enumerate}[leftmargin=1.2em]
\item O utilizador introduz contacto e PIN;
\item O sistema valida as credenciais;
\item O sistema abre o portal do contribuinte ou do administrador.
\end{enumerate}}
{A1: credenciais incorrectas: o sistema recusa e conta a tentativa. A2: após tentativas sucessivas falhadas, o acesso é bloqueado temporariamente.}

\ucspec{UC06}{Consulta de histórico e recibos (derivado)}
{Cidadão ou comerciante.}
{Permite consultar apenas os próprios pagamentos e reabrir cada recibo digital.}
{O utilizador está autenticado.}
{Lista de pagamentos e recibo apresentados.}
{\begin{enumerate}[leftmargin=1.2em]
\item O utilizador abre ``Meus pagamentos'';
\item O sistema lista os pagamentos do utilizador;
\item O utilizador seleciona um pagamento e visualiza o recibo.
\end{enumerate}}
{A1: sem pagamentos: o sistema indica que ainda não há pagamentos e sugere ``Pagar taxa''.}

\begin{table}[!htbp]
\caption{Casos de uso UC07 e UC08 (derivados)}\label{tab:UC08}
\begingroup\singlespacing\small
\begin{tabularx}{\textwidth}{@{}>{\bfseries}p{1.6cm}p{3.4cm}X@{}}
\toprule
ID & Nome / actor & Descrição resumida \\
\midrule
UC07 & Consultar auditoria (administrador) & Lista cronológica de acções (quem, quando, o quê): registos, pagamentos, alterações de taxas, entradas e saídas. \\
UC08 & Registar pagamento presencial (agente municipal) & O agente seleciona o contribuinte e a taxa, recebe numerário ou usa POS, e o sistema emite o recibo com a mesma lógica do UC02. \\
\bottomrule
\end{tabularx}
\endgroup
\end{table}

\begin{table}[!htbp]
\caption{Matriz de rastreabilidade entre requisitos funcionais e casos de uso}\label{tab:rastr}
\begingroup\singlespacing\small
\begin{tabular}{@{}lccccccccc@{}}
\toprule
 & RF01 & RF02 & RF03 & RF04 & RF05 & RF06 & RF07 & RF08 & RF09 \\
\midrule
UC01 & x & & & & & x & x & x & x \\
UC02 & & x & x & & & x & x & & x \\
UC03 & & & & x & & & x & & \\
UC04 & & & & & x & & x & & \\
UC05 & x & & & & & & x & & \\
UC06 & & & x & & & & & & \\
UC07 & & & & & & & x & & \\
UC08 & & x & x & & & x & x & & \\
\bottomrule
\end{tabular}
\endgroup
\end{table}

\section{Modelo de classes}
A Figura~\ref{fig:classes} apresenta o diagrama de classes de domínio. \texttt{Contribuinte} e \texttt{Administrador} especializam \texttt{Utilizador}; um contribuinte efectua vários pagamentos, cada um referente a uma taxa e gerando um recibo e notificações; as taxas pertencem a uma categoria e são geridas pelo administrador; todas as acções relevantes geram registos de auditoria.

\begin{figure}[!htbp]
\centering
\begin{adjustbox}{max width=\textwidth}
\begin{tikzpicture}[font=\scriptsize,
 cls/.style={rectangle split,rectangle split parts=3,rectangle split part align={center,left,left},draw,text width=4.2cm,inner sep=3pt},
 inh/.style={-{Triangle[open,length=3mm,width=3mm]}}]
\node[cls] (ut) at (0,0) {\textbf{Utilizador}\nodepart{two}id\\nome\\telefone\\morada\\pinHash\\papel\nodepart{three}autenticar()\\alterarPin()};
\node[cls] (cw) at (-5.5,0) {\textbf{CarteiraMovel}\nodepart{two}id\\operadora\\numero\nodepart{three}validar()};
\node[cls] (au) at (5.5,0) {\textbf{Auditoria}\nodepart{two}id\\accao\\entidade\\detalhe\\dataHora\nodepart{three}registar()};
\node[cls] (co) at (-5.5,-4.6) {\textbf{Contribuinte}\nodepart{two}tipo (cidadão/comerciante)\nodepart{three}registar()\\pagarTaxa()\\consultarHistorico()};
\node[cls] (ad) at (5.5,-4.6) {\textbf{Administrador}\nodepart{two}perfil (admin/gestor)\nodepart{three}gerirTaxas()\\gerarRelatorio()\\consultarAuditoria()};
\node[cls] (pa) at (-5.5,-9.2) {\textbf{Pagamento}\nodepart{two}referencia\\valorTotal\\quantidade\\metodo\\estado\\dataHora\nodepart{three}calcularValor()\\confirmar()\\falhar()};
\node[cls] (ta) at (0,-9.2) {\textbf{Taxa}\nodepart{two}id\\nome\\valor\\periodicidade\\activa\nodepart{three}criar()\\alterar()\\desactivar()};
\node[cls] (ca) at (5.5,-9.2) {\textbf{Categoria}\nodepart{two}id\\nome\nodepart{three}criar()\\renomear()};
\node[cls] (re) at (-5.5,-13.6) {\textbf{Recibo}\nodepart{two}numero\\emitidoEm\nodepart{three}gerar()\\exportarPdf()};
\node[cls] (no) at (0,-13.6) {\textbf{Notificacao}\nodepart{two}canal (SMS/e-mail)\\mensagem\\estado\nodepart{three}enviar()};
\draw[inh] (co.north) -- ++(0,.8) -| ([xshift=-.4cm]ut.south);
\draw[inh] (ad.north) -- ++(0,.8) -| ([xshift=.4cm]ut.south);
\draw (co) -- node[pos=.2,right]{1} node[pos=.8,right]{0..*} (cw);
\draw (co) -- node[pos=.2,right]{1} node[pos=.8,right]{0..*} (pa);
\draw (pa) -- node[pos=.15,above]{0..*} node[pos=.85,above]{1} (ta);
\draw (ta) -- node[pos=.15,above]{0..*} node[pos=.85,above]{1} (ca);
\draw (ad) -- node[pos=.2,right]{1} node[pos=.8,right]{0..*} (ca);
\draw (pa) -- node[pos=.2,right]{1} node[pos=.8,right]{1} (re);
\draw (pa) -- node[pos=.25,above]{1} node[pos=.75,above]{0..*} (no);
\draw (ut) -- node[pos=.15,above]{1} node[pos=.85,above]{0..*} (au);
\end{tikzpicture}
\end{adjustbox}
\caption{Diagrama de classes de domínio do SIIFM.}\label{fig:classes}
\end{figure}

\section{Diagramas de sequência}
\label{sec:seq}
\newcommand{\lifeline}[4]{\node[draw,fill=gray!10,minimum width=2.1cm,minimum height=.7cm,font=\scriptsize,align=center] at (#1,0) {#2};\draw[dashed,gray] (#1,-.35) -- (#1,-#3);}
\newcommand{\msg}[5][solid]{\draw[-{Stealth[length=2mm]},#1] (#2,#4) -- node[above,midway,font=\scriptsize,inner sep=1pt]{#5} (#3,#4);}
A Figura~\ref{fig:seqpag} descreve o fluxo de pagamento (UC02) e a Figura~\ref{fig:seqreg} descreve o registo de contribuinte (UC01).

\begin{figure}[!htbp]
\centering
\begin{adjustbox}{max width=\textwidth}
\begin{tikzpicture}
\lifeline{0}{Contribuinte}{13.4}{}
\lifeline{2.8}{Frontend\\(PWA)}{13.4}{}
\lifeline{5.6}{API\\Backend}{13.4}{}
\lifeline{8.4}{Operadora\\móvel}{13.4}{}
\lifeline{11.2}{Base de\\dados}{13.4}{}
\lifeline{14}{Serviço\\SMS/e-mail}{13.4}{}
\msg{0}{2.8}{-1.0}{1: Selecciona taxa}
\msg{2.8}{5.6}{-1.8}{2: GET /taxas/\{id\}}
\msg{5.6}{11.2}{-2.6}{3: Consulta taxa}
\msg[dashed]{11.2}{5.6}{-3.4}{4: Taxa e valor}
\draw[-{Stealth[length=2mm]}] (5.6,-4.2) -- (6.1,-4.2) -- (6.1,-4.7) -- (5.6,-4.7);
\node[right,font=\scriptsize,text width=2.2cm] at (6.15,-4.45) {5: Calcula valor e gera referência};
\msg[dashed]{5.6}{2.8}{-5.6}{6: Valor e referência}
\msg[dashed]{2.8}{0}{-6.4}{7: Mostra referência}
\msg{0}{2.8}{-7.2}{8: Confirma pagamento}
\msg{2.8}{5.6}{-8.0}{9: POST /pagamentos}
\msg{5.6}{8.4}{-8.8}{10: Pede débito}
\msg[dashed]{8.4}{5.6}{-9.6}{11: Autorizado}
\msg{5.6}{11.2}{-10.4}{12: Grava pagamento, recibo e auditoria}
\msg{5.6}{14}{-11.2}{13: Envia recibo (SMS/e-mail)}
\msg[dashed]{5.6}{2.8}{-12.0}{14: Recibo digital}
\msg[dashed]{2.8}{0}{-12.8}{15: Apresenta recibo}
\end{tikzpicture}
\end{adjustbox}
\caption{Diagrama de sequência do pagamento de uma taxa (UC02).}\label{fig:seqpag}
\end{figure}

\begin{figure}[!htbp]
\centering
\begin{adjustbox}{max width=\textwidth}
\begin{tikzpicture}
\lifeline{0}{Utilizador}{8.2}{}
\lifeline{3}{Frontend}{8.2}{}
\lifeline{6}{API\\Backend}{8.2}{}
\lifeline{9}{Base de\\dados}{8.2}{}
\lifeline{12}{Serviço\\SMS/e-mail}{8.2}{}
\msg{0}{3}{-1.0}{1: Preenche dados}
\msg{3}{6}{-1.8}{2: POST /utilizadores}
\draw[-{Stealth[length=2mm]}] (6,-2.6) -- (6.5,-2.6) -- (6.5,-3.1) -- (6,-3.1);
\node[right,font=\scriptsize] at (6.55,-2.85) {3: Valida dados};
\msg{6}{9}{-3.9}{4: Contacto já existe?}
\msg[dashed]{9}{6}{-4.7}{5: Resposta}
\msg{6}{9}{-5.5}{6: Cria perfil (PIN em hash)}
\msg{6}{12}{-6.3}{7: Envia SMS de confirmação}
\msg[dashed]{6}{3}{-7.1}{8: 201 Criado}
\msg[dashed]{3}{0}{-7.9}{9: Mostra confirmação}
\end{tikzpicture}
\end{adjustbox}
\caption{Diagrama de sequência do registo de contribuinte (UC01). Se o contacto já existir (passo 5), a API devolve 409, o passo 6 não ocorre e o \emph{frontend} alerta duplicação.}\label{fig:seqreg}
\end{figure}

\section{Diagrama de actividades}
A Figura~\ref{fig:act} modela o fluxo de actividades do pagamento, incluindo a ramificação por método e o caso de falha.

\begin{figure}[!htbp]
\centering
\begin{adjustbox}{max width=\textwidth}
\begin{tikzpicture}[font=\small,
 act/.style={draw,rounded corners=6pt,minimum width=4.6cm,minimum height=.8cm,align=center},
 dec/.style={diamond,draw,aspect=2.3,inner sep=1pt,font=\scriptsize,align=center},
 ar/.style={-{Stealth[length=2mm]}}]
\node[circle,fill=black,minimum size=.35cm] (s) at (0,0) {};
\node[act] (a1) at (0,-1.3) {Autenticar-se};
\node[act] (a2) at (0,-2.6) {Seleccionar taxa e períodos};
\node[act] (a3) at (0,-3.9) {Calcular valor e gerar referência};
\node[act] (a4) at (0,-5.2) {Escolher método de pagamento};
\node[dec] (d1) at (0,-6.7) {Carteira\\móvel?};
\node[act,minimum width=3.6cm] (a5) at (-2.8,-8.3) {Autorizar na operadora};
\node[act,minimum width=3.6cm] (a6) at (2.8,-8.3) {Agente regista\\POS / numerário};
\node[dec] (d2) at (0,-10.1) {Confirmado?};
\node[act,minimum width=2.6cm] (a9) at (-5.4,-10.1) {Informar falha};
\node[act] (a7) at (0,-11.8) {Gravar pagamento e auditoria};
\node[act] (a8) at (0,-13.1) {Emitir recibo e notificar};
\node[circle,draw,minimum size=.4cm,inner sep=0] (e) at (0,-14.4) {\tikz\fill (0,0) circle (.12);};
\draw[ar] (s)--(a1); \draw[ar] (a1)--(a2); \draw[ar] (a2)--(a3); \draw[ar] (a3)--(a4); \draw[ar] (a4)--(d1);
\draw[ar] (d1.west) -| node[pos=.2,above]{sim} (a5.north);
\draw[ar] (d1.east) -| node[pos=.2,above]{não} (a6.north);
\draw (a5.south) -- ++(0,-.4) -| (d2.north);
\draw (a6.south) -- ++(0,-.4) -| (d2.north);
\draw[ar] (d2) -- node[right]{sim} (a7);
\draw[ar] (d2.west) -- node[above]{não} (a9.east);
\draw[ar] (a9.north) |- (a4.west);
\draw[ar] (a7)--(a8); \draw[ar] (a8)--(e);
\end{tikzpicture}
\end{adjustbox}
\caption{Diagrama de actividades do pagamento de taxas.}\label{fig:act}
\end{figure}

\section{Modelo de dados}
O modelo entidade-relacionamento (Figura~\ref{fig:er}) é normalizado até à terceira forma normal. O dicionário de dados completo consta do Apêndice~A. As chaves estrangeiras garantem a integridade referencial; as taxas não são apagadas fisicamente (campo \texttt{activa}) para preservar o histórico de pagamentos.

\begin{figure}[!htbp]
\centering
\begin{adjustbox}{max width=\textwidth}
\begin{tikzpicture}[font=\scriptsize,
 ent/.style={rectangle split,rectangle split parts=2,rectangle split part align={center,left},draw,text width=4.2cm,inner sep=3pt}]
\node[ent] (ut) at (0,0) {\textbf{UTILIZADOR}\nodepart{two}\underline{id} (PK)\\nome\\telefone (único)\\morada\\tipo\\papel\\pin\_hash\\criado\_em};
\node[ent] (cw) at (-5.5,0) {\textbf{CARTEIRA\_MOVEL}\nodepart{two}\underline{id} (PK)\\utilizador\_id (FK)\\operadora\\numero};
\node[ent] (au) at (5.5,0) {\textbf{AUDITORIA}\nodepart{two}\underline{id} (PK)\\utilizador\_id (FK)\\accao\\entidade\\detalhe\\data\_hora};
\node[ent] (re) at (-5.5,-5.4) {\textbf{RECIBO}\nodepart{two}\underline{id} (PK)\\pagamento\_id (FK, único)\\numero (único)\\emitido\_em};
\node[ent] (pa) at (0,-5.8) {\textbf{PAGAMENTO}\nodepart{two}\underline{id} (PK)\\referencia (única)\\utilizador\_id (FK)\\taxa\_id (FK)\\carteira\_id (FK, opcional)\\metodo\\quantidade\\valor\_total\\estado\\data\_hora};
\node[ent] (ta) at (5.5,-5.8) {\textbf{TAXA}\nodepart{two}\underline{id} (PK)\\categoria\_id (FK)\\nome\\valor\\periodicidade\\activa};
\node[ent] (no) at (0,-11.4) {\textbf{NOTIFICACAO}\nodepart{two}\underline{id} (PK)\\pagamento\_id (FK)\\canal\\mensagem\\estado\\enviada\_em};
\node[ent] (ca) at (5.5,-11.4) {\textbf{CATEGORIA}\nodepart{two}\underline{id} (PK)\\nome (único)};
\draw (ut) -- node[pos=.15,above]{1} node[pos=.85,above]{N} (cw);
\draw (ut) -- node[pos=.15,above]{1} node[pos=.85,above]{N} (au);
\draw (ut) -- node[pos=.12,right]{1} node[pos=.88,right]{N} (pa);
\draw (cw) -- node[pos=.15,left]{1} node[pos=.85,left]{N} (pa);
\draw (pa) -- node[pos=.15,above]{1} node[pos=.85,above]{1} (re);
\draw (ta) -- node[pos=.15,above]{1} node[pos=.85,above]{N} (pa);
\draw (ca) -- node[pos=.15,right]{1} node[pos=.85,right]{N} (ta);
\draw (pa) -- node[pos=.15,right]{1} node[pos=.85,right]{N} (no);
\end{tikzpicture}
\end{adjustbox}
\caption{Modelo entidade-relacionamento do SIIFM.}\label{fig:er}
\end{figure}

\section{Arquitectura técnica}
O SIIFM adopta uma arquitectura em camadas (Figura~\ref{fig:arq}): apresentação (três portais React separados: cidadão/comerciante com PWA, agente de balcão e administração), aplicação (API RESTful em Node.js/Express ou Python/Django REST Framework), dados (PostgreSQL ou MySQL) e integrações externas (operadoras móveis e fornecedor de SMS). A infraestrutura pode ser em nuvem (AWS ou Azure) ou \emph{on-premises}, com cópias de segurança automáticas, \emph{logs} auditáveis e TLS em trânsito.

\begin{figure}[!htbp]
\centering
\begin{adjustbox}{max width=\textwidth}
\begin{tikzpicture}[font=\small,
 bx/.style={draw,rounded corners=3pt,minimum height=1cm,align=center,text width=3.6cm,fill=gray!8},
 ar/.style={-{Stealth[length=2mm]}}]
\node[font=\scriptsize\itshape,fill=white,inner sep=1pt] at (7,1.2) {Camada de apresentação};
\node[bx] (p1) at (1.9,0) {Portal do cidadão\\React + PWA};
\node[bx] (p2) at (7,0) {Portal de administração\\(React)};
\node[bx] (p3) at (12.1,0) {Portal do agente\\de balcão (React)};
\draw[rounded corners] (-.2,-2) rectangle (14.2,-6.2);
\node[font=\scriptsize\itshape,fill=white,inner sep=1pt] at (7,-1.7) {Camada de aplicação: API RESTful sobre HTTPS (TLS 1.3)};
\node[bx,minimum height=.8cm,text width=2.8cm] (m1) at (1.7,-3) {Autenticação e RBAC};
\node[bx,minimum height=.8cm,text width=2.8cm] (m2) at (5.2,-3) {Contribuintes};
\node[bx,minimum height=.8cm,text width=2.8cm] (m3) at (8.7,-3) {Taxas e categorias};
\node[bx,minimum height=.8cm,text width=2.8cm] (m4) at (12.2,-3) {Pagamentos};
\node[bx,minimum height=.8cm,text width=2.8cm] (m5) at (1.7,-4.9) {Recibos};
\node[bx,minimum height=.8cm,text width=2.8cm] (m6) at (5.2,-4.9) {Relatórios};
\node[bx,minimum height=.8cm,text width=2.8cm] (m7) at (8.7,-4.9) {Notificações};
\node[bx,minimum height=.8cm,text width=2.8cm] (m8) at (12.2,-4.9) {Auditoria e logs};
\node[font=\scriptsize\itshape,fill=white,inner sep=1pt] at (7,-7.2) {Camada de dados e integrações externas};
\node[bx,text width=2.9cm] (d1) at (1.9,-8.6) {PostgreSQL / MySQL};
\node[bx,text width=2.9cm] (d2) at (5.4,-8.6) {Cópias de segurança\\e armazenamento de logs};
\node[bx,text width=2.9cm] (x1) at (8.9,-8.6) {Operadoras móveis\\(M-Pesa, e-Mola)};
\node[bx,text width=2.9cm] (x2) at (12.4,-8.6) {Fornecedor de\\SMS};
\draw[ar] (p1)--(1.9,-2); \draw[ar] (p2)--(7,-2); \draw[ar] (p3)--(12.1,-2);
\draw[ar] (1.9,-6.2)--(d1); \draw[ar] (5.4,-6.2)--(d2);
\draw[ar] (8.9,-6.2)--(x1); \draw[ar] (12.4,-6.2)--(x2);
\end{tikzpicture}
\end{adjustbox}
\caption{Arquitectura em camadas do SIIFM.}\label{fig:arq}
\end{figure}

\section{Segurança e protecção de dados}
\begin{table}[!htbp]
\caption{Medidas de segurança e conformidade}\label{tab:seg}
\begingroup\singlespacing\small
\begin{tabularx}{\textwidth}{@{}>{\bfseries}p{3.8cm}X@{}}
\toprule
Ameaça / exigência & Medida prevista \\
\midrule
Intercepção de dados & TLS~1.3 em todas as comunicações; \emph{HSTS}. \\
Roubo de credenciais & PIN e palavras-passe guardados com \emph{hash} e \emph{salt} (por exemplo, bcrypt ou Argon2); bloqueio após tentativas falhadas. \\
Acesso indevido & Controlo de acessos por perfil (RBAC); o contribuinte só acede aos seus próprios dados. \\
Injecção e XSS & Consultas parametrizadas, validação e escape de entradas, \emph{Content-Security-Policy}. \\
Falsificação de confirmações & \emph{Callbacks} da operadora com assinatura (HMAC) e verificação de \emph{idempotência} para evitar duplo registo. \\
Fraude interna & Registo de auditoria imutável de todas as transacções e alterações de taxas. \\
Perda de dados & Cópias automáticas, teste periódico de restauro. \\
Protecção de dados pessoais & Minimização de dados, finalidade declarada, encriptação de dados sensíveis em repouso e revisão jurídica (RGPD e lei moçambicana). \\
\bottomrule
\end{tabularx}
\endgroup
\end{table}

\section{Plano de testes}
A Tabela~\ref{tab:testes} indica o estado de cada teste. Os testes marcados como \emph{Executado} correspondem a 39 testes automáticos do \emph{backend} (comando \texttt{npm test}) e a 7 testes da interface React (\texttt{npm run test:client}), todos aprovados. Os testes das integrações usam servidores falsos que imitam a Twilio, a Africa's Talking e a M-Pesa: verificam o formato dos pedidos e o tratamento das respostas, mas não provam que os serviços reais aceitam esses pedidos. O teste T09 está \emph{Parcial}: foram automatizadas a rejeição de tokens adulterados e a tentativa de injecção SQL, mas não houve teste de intrusão externo. Os testes de carga, de recuperação e de usabilidade (T10 a T12) ainda não foram executados.

\begin{table}[!htbp]
\caption{Plano de testes}\label{tab:testes}
\begingroup\singlespacing\small
\begin{tabularx}{\textwidth}{@{}>{\bfseries}p{1cm}p{1.9cm}Xp{3cm}p{1.8cm}@{}}
\toprule
ID & Tipo & Caso de teste & Resultado esperado & Estado \\
\midrule
T01 & Unitário & Validar telemóvel (82--87 + 7 dígitos) e PIN (mín. 4 dígitos). & Entradas inválidas rejeitadas.  & Executado \\
T02 & Unitário & Calcular valor total = valor da taxa $\times$ períodos. & Valor correcto.  & Executado \\
T03 & Integração & Registo grava utilizador e dispara SMS. & Perfil criado e SMS enviado.  & Executado \\
T04 & Integração & Pagamento confirmado grava pagamento, recibo e auditoria na mesma transacção. & Os três registos existem; nenhum parcial.  & Executado \\
T05 & Integração & \emph{Callback} duplicado da operadora. & Apenas um pagamento registado.  & Executado \\
T06 & E2E & Registo $\to$ login $\to$ pagamento $\to$ recibo $\to$ relatório do administrador. & Valor aparece no relatório.  & Executado \\
T07 & E2E & Contacto duplicado no registo. & Alerta de duplicação.  & Executado \\
T08 & Segurança & Contribuinte tenta aceder a rotas do administrador e a recibos de outro. & Acesso negado (403).  & Executado \\
T09 & Segurança & Tentativas de injecção SQL/XSS; teste de intrusão. & Sem vulnerabilidades críticas.  & Parcial \\
T10 & Desempenho & Carga com 50\,000 utilizadores registados e pico de pagamentos. & Latência $<$ 2~s nas acções críticas.  & Por executar \\
T11 & Disponibilidade & Simulação de falha e restauro. & Serviço recuperado; sem perda de pagamentos confirmados.  & Por executar \\
T12 & Usabilidade & Utilizadores reais executam registo e pagamento; verificação WCAG básica. & Tarefas concluídas sem ajuda; contraste e teclado adequados.  & Por executar \\
\bottomrule
\end{tabularx}
\endgroup
\end{table}

\section{Plano de desenvolvimento (Scrum)}
\label{sec:scrum}
O desenvolvimento seguiria \emph{sprints} de duas a três semanas, com \emph{Product Backlog} priorizado por valor de negócio e demonstrações aos \emph{stakeholders} municipais no final de cada \emph{sprint}.

\begin{table}[!htbp]
\caption{Plano de \emph{sprints} (estimativa)}\label{tab:sprints}
\begingroup\singlespacing\small
\begin{tabularx}{\textwidth}{@{}>{\bfseries}p{2.3cm}p{2.2cm}X@{}}
\toprule
\emph{Sprint} & Duração & Objectivo e entregas \\
\midrule
0 & 1 semana & Preparação: ambiente, repositório, \emph{pipeline} CI/CD, desenho de base de dados. \\
1 & 2 semanas & Autenticação (RBAC) e registo de contribuintes (UC01, UC05). \\
2 & 2 semanas & Gestão de taxas e categorias (UC03); auditoria base (UC07). \\
3 & 3 semanas & Pagamentos e integração com M-Pesa/e-Mola em \emph{sandbox} (UC02, UC08). \\
4 & 2 semanas & Recibos digitais, histórico e notificações SMS/e-mail (UC06). \\
5 & 3 semanas & Relatórios, \emph{dashboards} e exportação PDF/Excel (UC04); canal USSD. \\
6 & 3 semanas & Testes de segurança e de desempenho, correcções, \emph{deploy} e aceitação (UAT). \\
\bottomrule
\end{tabularx}
\endgroup
\end{table}

\section{Gestão de riscos}
\begin{table}[!htbp]
\caption{Principais riscos e mitigação}\label{tab:riscos}
\begingroup\singlespacing\small
\begin{tabularx}{\textwidth}{@{}p{3.4cm}p{1.6cm}X@{}}
\toprule
Risco & Impacto & Mitigação \\
\midrule
Atraso ou condições comerciais na integração com operadoras & Alto & Iniciar contactos na \emph{Sprint}~0; usar \emph{sandbox}; camada de abstracção por operadora. \\
Conectividade limitada & Médio & PWA com \emph{offline} parcial; canais USSD/SMS. \\
Baixa adopção pelos contribuintes & Alto & Pagamento presencial por POS/numerário; formação e campanhas. \\
Falhas na protecção de dados & Alto & Segurança desde o desenho; revisão jurídica; testes de intrusão. \\
Indisponibilidade do serviço & Médio & Redundância, \emph{backups}, monitorização (meta de 99\%). \\
Fraude interna & Alto & Auditoria imutável e segregação de funções. \\
\bottomrule
\end{tabularx}
\endgroup
\end{table}

\section{Sistema implementado}
\label{sec:proto}
Além da modelagem, o SIIFM foi implementado como aplicação completa, cujo código é entregue com este relatório. A Tabela~\ref{tab:impl} resume a implementação e a sua relação com a arquitectura da Figura~\ref{fig:arq}.

\begin{table}[!htbp]
\caption{Componentes do sistema implementado}\label{tab:impl}
\begingroup\singlespacing\small
\begin{tabularx}{\textwidth}{@{}>{\bfseries}p{3cm}X@{}}
\toprule
Componente & Implementação \\
\midrule
Backend & Node.js 22 e Express; API REST com os \emph{endpoints} do Apêndice~B; regras de negócio (cálculo do valor, estados do pagamento, recibos) e transacções atómicas. \\
Base de dados & SQLite (módulo \texttt{node:sqlite}) com o esquema do modelo ER, chaves estrangeiras, restrições \texttt{CHECK} e gatilhos que tornam a tabela de auditoria imutável. A migração para PostgreSQL ou MySQL está prevista na arquitectura. \\
Segurança & Modo de produção que recusa arrancar com segredos fracos ou fornecedores simulados; PIN com \emph{scrypt} e \emph{salt}; sessões por \emph{token} assinado (HMAC); controlo de acessos por perfil; bloqueio após cinco tentativas falhadas; cabeçalhos de segurança e \emph{Content-Security-Policy}; \emph{callbacks} da operadora assinados (HMAC-SHA256) e idempotentes, com verificação do valor. \\
Operadora móvel & \textbf{M-Pesa:} adaptador da API M-Pesa de Moçambique (Vodacom, pagamento C2B), com recusa de códigos desconhecidos como sucesso e reconciliação periódica por consulta de estado. Foi testado contra um servidor falso e não contra o serviço real. \textbf{e-Mola:} apenas simulador (sem API pública documentada); em produção fica indisponível. \\
Interface web & React~18 (Vite), organizada em componentes e um ficheiro por ecrã, com \textbf{três portais separados}, cada um com página, \emph{login} e sessão próprios: cidadão/comerciante (\texttt{/}, PWA instalável), agente de balcão (\texttt{/agente/}) e administração (\texttt{/admin/}, administrador e gestor). A API recusa o \emph{login} de uma conta no portal errado e continua a proteger cada rota por perfil. Um aviso visível indica o que está simulado. \\
Notificações & SMS por adaptadores HTTP para Twilio e Africa's Talking, enviados em segundo plano (uma falha de SMS nunca bloqueia um pagamento); o estado de cada SMS fica registado. Em demonstração usa-se um simulador. \\
Testes & 39 testes automáticos do \emph{backend} (API, adaptadores de SMS e M-Pesa contra servidores falsos, portais e modo de produção) e 7 testes da interface React (React Testing Library e Vitest) executados contra o servidor real (Tabela~\ref{tab:testes}). \\
\bottomrule
\end{tabularx}
\endgroup
\end{table}

O pagamento por carteira móvel segue o fluxo da Figura~\ref{fig:seqpag}: o contribuinte escolhe a taxa e o método, o sistema gera a referência e pede o débito à operadora; a confirmação chega por \emph{callback}, o sistema grava de forma atómica o pagamento, o recibo, a notificação e a auditoria, e a interface apresenta o recibo. O pagamento em POS ou numerário gera uma referência que o agente confirma no balcão. Uma demonstração estática da interface (dados guardados apenas no navegador, sem servidor) está disponível em \url{https://claude.ai/artifact/4L2dSkePcy3LHdfxnJyAVn}. As instruções de instalação constam do Apêndice~C.

\section{Discussão das hipóteses}
A modelagem permite uma avaliação \emph{conceptual} das hipóteses; a confirmação empírica exige a implementação e a utilização real.

\begin{table}[!htbp]
\caption{Discussão das hipóteses}\label{tab:hip}
\begingroup\singlespacing\small
\begin{tabularx}{\textwidth}{@{}>{\bfseries}p{1cm}X@{}}
\toprule
H & Avaliação com base na modelagem \\
\midrule
H1 & Suportada em desenho: UC02 e UC08 unificam carteira móvel, POS e numerário, com cálculo automático e recibo imediato. O ganho efectivo em tempo e perdas só pode ser medido em piloto. \\
H2 & Suportada em desenho: perfil único (UC01, RF08), histórico (UC06) e auditoria (UC07, RF07). Demonstrada no sistema implementado. \\
H3 & Parcialmente suportada: a PWA e o canal USSD/SMS estão previstos (RF09); o USSD não foi implementado no protótipo. \\
H4 & Suportada em desenho e demonstrada no sistema implementado com totais por taxa, método e dia; a actualização em tempo real depende do \emph{backend}. \\
\bottomrule
\end{tabularx}
\endgroup
\end{table}
%% =====================================================
\chapter{Conclusões, Limitações e Recomendações}

\section{Conclusões}
O trabalho cumpriu o objectivo geral de modelar e conceber o SIIFM. Foram identificados sete actores, nove requisitos funcionais e sete não funcionais, e modelados oito casos de uso, o diagrama de classes, dois diagramas de sequência, um diagrama de actividades, o modelo entidade-relacionamento e a arquitectura em camadas. Foram ainda definidos o plano de testes, o plano de \emph{sprints}, as medidas de segurança e a análise de riscos, e implementado o sistema (API, base de dados e interface com portais por perfil), verificado por testes automáticos.

Quanto às hipóteses, a modelagem sustenta H1, H2 e H4 ao nível do desenho, e H3 parcialmente (o canal USSD está especificado mas não foi implementado). A escolha de uma arquitectura multicanal (carteira móvel, POS e numerário) responde ao facto de nem todos os contribuintes possuírem carteira móvel, e a auditoria de todas as operações responde ao problema de rastreabilidade identificado na problematização.

\section{Limitações}
\begin{itemize}
\item \textbf{Sem validação empírica:} o modelo não foi testado com contribuintes nem com funcionários municipais; as hipóteses foram avaliadas apenas ao nível conceptual;
\item \textbf{Integrações reais não validadas:} os adaptadores de SMS (Twilio, Africa's Talking) e de M-Pesa foram escritos a partir da documentação pública e testados apenas contra servidores falsos, por não haver credenciais; têm de ser validados em \emph{sandbox} antes de usar dinheiro real. O e-Mola só existe como simulador, por falta de API documentada, e o USSD não foi implementado. As condições comerciais e regulamentares das operadoras não foram verificadas;
\item \textbf{Implementação de demonstração:} usa SQLite em vez de PostgreSQL ou MySQL, não termina TLS por si (exige um \emph{reverse proxy} HTTPS), não tem limitação global de pedidos nem recuperação de PIN, e a restrição de rede aos portais de administração e do agente fica por configurar e não gera PDF nativo (usa a impressão do navegador); a demonstração estática publicada guarda dados apenas no navegador. Nenhuma das versões deve ser usada com dados reais sem as correcções de produção;
\item \textbf{Aspectos jurídicos:} a conformidade com a legislação moçambicana de protecção de dados e com os regulamentos de pagamentos não foi validada por jurista;
\item \textbf{Estimativas:} as durações das \emph{sprints} e os requisitos de desempenho (50\,000 utilizadores, 2~s) não foram medidos;
\item \textbf{Âmbito:} não foram detalhados aspectos como multas e juros, isenções, reembolsos, conciliação financeira e integração com o sistema contabilístico do município.
\end{itemize}

\section{Recomendações e trabalho futuro}
\begin{enumerate}
\item Migrar para PostgreSQL ou MySQL, colocar o sistema atrás de HTTPS e executar os testes que faltam (T09 a T12 da Tabela~\ref{tab:testes});
\item Validar o adaptador M-Pesa e o SMS em \emph{sandbox} (contrato de comerciante com a Vodacom e conta no fornecedor de SMS), implementar o e-Mola com a documentação da Movitel e só depois passar ao modo de produção;
\item Realizar um piloto num mercado ou bairro, medindo o tempo de pagamento, o valor arrecadado e a satisfação;
\item Rever o desenho com um jurista antes de qualquer tratamento de dados reais;
\item Acrescentar funcionalidades de multas, isenções, reembolsos e conciliação financeira.
\end{enumerate}

%% =====================================================
\begin{thebibliography}{99}
\addcontentsline{toc}{chapter}{Bibliografia}
\renewcommand{\bibname}{Bibliografia}
\bibitem{booch2005} BOOCH, G.; RUMBAUGH, J.; JACOBSON, I. \emph{The Unified Modeling Language User Guide}. 2.ª ed. Addison-Wesley, 2005.
\bibitem{elmasri2016} ELMASRI, R.; NAVATHE, S. B. \emph{Fundamentals of Database Systems}. 7.ª ed. Pearson, 2016.
\bibitem{enif} REPÚBLICA DE MOÇAMBIQUE; BANCO DE MOÇAMBIQUE. \emph{Estratégia Nacional de Inclusão Financeira 2016--2022}. Maputo.
\bibitem{fielding2000} FIELDING, R. T. \emph{Architectural Styles and the Design of Network-based Software Architectures}. Tese de doutoramento, University of California, Irvine, 2000.
\bibitem{findex2021} DEMIRGÜÇ-KUNT, A.; KLAPPER, L.; SINGER, D.; ANSAR, S. \emph{The Global Findex Database 2021}. Washington, D.C.: World Bank, 2022.
\bibitem{fowler2004} FOWLER, M. \emph{UML Distilled}. 3.ª ed. Addison-Wesley, 2004.
\bibitem{jack2011} JACK, W.; SURI, T. \emph{Mobile Money: The Economics of M-PESA}. NBER Working Paper n.º 16721, 2011.
\bibitem{larman2004} LARMAN, C. \emph{Applying UML and Patterns}. 3.ª ed. Prentice Hall, 2004.
\bibitem{lei2017} REPÚBLICA DE MOÇAMBIQUE. \emph{Lei n.º 3/2017, de 9 de Janeiro (Lei das Transacções Electrónicas)}. Boletim da República.
\bibitem{myers2011} MYERS, G. J.; SANDLER, C.; BADGETT, T. \emph{The Art of Software Testing}. 3.ª ed. Wiley, 2011.
\bibitem{nielsen1993} NIELSEN, J. \emph{Usability Engineering}. Academic Press, 1993.
\bibitem{owasp2021} OWASP FOUNDATION. \emph{OWASP Top 10 -- 2021}. Disponível em: \url{https://owasp.org/Top10/}.
\bibitem{pressman2020} PRESSMAN, R. S.; MAXIM, B. R. \emph{Software Engineering: A Practitioner's Approach}. 9.ª ed. McGraw-Hill, 2020.
\bibitem{rgpd2016} UNIÃO EUROPEIA. \emph{Regulamento (UE) 2016/679 (Regulamento Geral sobre a Protecção de Dados)}. Jornal Oficial da União Europeia, 2016.
\bibitem{scrum2020} SCHWABER, K.; SUTHERLAND, J. \emph{The Scrum Guide}. 2020. Disponível em: \url{https://scrumguides.org}.
\bibitem{sommerville2016} SOMMERVILLE, I. \emph{Software Engineering}. 10.ª ed. Pearson, 2016.
\bibitem{wcag2018} W3C. \emph{Web Content Accessibility Guidelines (WCAG) 2.1}. 2018. Disponível em: \url{https://www.w3.org/TR/WCAG21/}.
\end{thebibliography}

%% =====================================================
\appendix
\renewcommand{\thechapter}{\Alph{chapter}}
\titleformat{\chapter}[hang]{\fontsize{14}{17}\selectfont\bfseries}{\MakeUppercase{Apêndice \Alph{chapter}} --}{.5em}{\MakeUppercase}

\chapter{Dicionário de dados}
\begin{table}[!htbp]
\caption{Dicionário de dados resumido}\label{tab:dd}
\begingroup\singlespacing\small
\begin{tabularx}{\textwidth}{@{}>{\ttfamily}p{3cm}X@{}}
\toprule
\textrm{Tabela} & Campos principais \\
\midrule
utilizador & id (PK), nome, telefone (único), morada, tipo (cidadão/comerciante), papel (contribuinte/administrador/gestor), pin\_hash, criado\_em. \\
carteira\_movel & id (PK), utilizador\_id (FK), operadora (M-Pesa/e-Mola), numero. \\
categoria & id (PK), nome (único). \\
taxa & id (PK), categoria\_id (FK), nome, valor (decimal), periodicidade (diária/mensal/anual/única), activa (booleano). \\
pagamento & id (PK), referencia (única), utilizador\_id (FK), taxa\_id (FK), carteira\_id (FK, opcional), metodo, quantidade, valor\_total, estado (pendente/confirmado/falhado), data\_hora. \\
recibo & id (PK), pagamento\_id (FK, único), numero (único), emitido\_em. \\
notificacao & id (PK), pagamento\_id (FK), canal (SMS/e-mail), mensagem, estado, enviada\_em. \\
auditoria & id (PK), utilizador\_id (FK), accao, entidade, detalhe, data\_hora. \\
\bottomrule
\end{tabularx}
\endgroup
\end{table}

\chapter{Interface da API REST}
\begin{table}[!htbp]
\caption{Endpoints principais da API}\label{tab:api}
\begingroup\singlespacing\small
\begin{tabularx}{\textwidth}{@{}>{\ttfamily}p{1.6cm}>{\ttfamily}p{5cm}X@{}}
\toprule
\textrm{Método} & \textrm{Rota} & Descrição e perfil \\
\midrule
POST & /auth/login & Autenticação (todos). \\
POST & /utilizadores & Registo de contribuinte (público); devolve 409 se duplicado. \\
GET & /utilizadores/me & Perfil do utilizador autenticado. \\
GET & /taxas & Lista de taxas activas (autenticado). \\
POST & /taxas & Criar taxa (administrador). \\
PUT & /taxas/\{id\} & Alterar taxa (administrador). \\
DELETE & /taxas/\{id\} & Desactivar taxa (administrador). \\
POST & /pagamentos & Iniciar pagamento e gerar referência (contribuinte/agente). \\
POST & /pagamentos/callback & Confirmação da operadora (assinada, idempotente). \\
GET & /pagamentos & Histórico (próprio; administrador vê todos). \\
GET & /recibos/\{numero\} & Recibo digital. \\
GET & /relatorios & Relatório por período/tipo; exporta PDF ou Excel (administrador/gestor). \\
GET & /auditoria & Registo de auditoria (administrador). \\
\bottomrule
\end{tabularx}
\endgroup
\end{table}

\chapter{Instalação e utilização do sistema}
\section*{Instalação}
Requisitos: Node.js 22.5 ou superior.
\begin{enumerate}
\item Extrair o projecto e, na pasta \texttt{siifm}, executar \texttt{npm install}. Este comando instala o \emph{backend} e compila a interface React (se a pasta \texttt{client/dist} ainda não existir);
\item Copiar \texttt{.env.example} para \texttt{.env} e definir \texttt{TOKEN\_SECRET} e \texttt{ADMIN\_PASSWORD}. Para SMS reais, \texttt{SMS\_PROVIDER} e as credenciais do fornecedor; para M-Pesa real, \texttt{MPESA\_MODE} e as credenciais da Vodacom; para operar com dinheiro real, \texttt{AMBIENTE=producao};
\item Executar \texttt{npm start} e abrir \texttt{http://localhost:3000} (portal do cidadão), \texttt{/admin/} (administração) e \texttt{/agente/} (agente de balcão);
\item Testes: \texttt{npm test} (\emph{backend}) e \texttt{npm run test:client} (interface React);
\item Após alterar a interface em \texttt{client/src}: \texttt{npm run build}. Em desenvolvimento, \texttt{npm run dev:api} e \texttt{npm run dev:client} (interface em \texttt{http://localhost:5173}).
\end{enumerate}
\section*{Utilização}
\begin{enumerate}
\item \textbf{Contribuinte:} em ``Registar-me'', indicar nome, telemóvel (82--87 e 7 dígitos), morada, tipo e PIN de 4 a 8 dígitos; entrar; em ``Pagar taxa'', escolher a taxa, o método e os períodos. Para M-Pesa, o número deve começar por 84 ou 85, e para e-Mola por 86 ou 87. No simulador, números terminados em 0 são recusados (saldo insuficiente).
\item \textbf{Administrador} (portal \texttt{/admin/}): entrar com o utilizador \texttt{admin} e a palavra-passe definida em \texttt{ADMIN\_PASSWORD} (por defeito \texttt{admin123}, a alterar); criar agentes e gestores em ``Utilizadores''; gerir taxas; consultar relatórios e auditoria.
\item \textbf{Agente de balcão} (portal \texttt{/agente/}): em ``Balcão'', confirmar o recebimento de referências geradas pelos contribuintes ou registar um pagamento presencial.
\end{enumerate}
\emph{Nota:} a palavra-passe por defeito e os simuladores de SMS e de operadora destinam-se apenas a demonstração; os portais mostram um aviso enquanto algum canal for simulado.

% =====================================================
% NOVO APÊNDICE D: EVIDÊNCIAS DA IMPLEMENTAÇÃO
% =====================================================
\chapter{Evidências da Implementação (Capturas de Ecrã)}
Este apêndice reúne as capturas de ecrã que evidenciam o funcionamento do sistema SIIFM implementado, comprovando os casos de uso e os requisitos descritos no Capítulo III.

\section*{Portal do Contribuinte}
\begin{figure}[!htbp]
\centering
\includegraphics[width=0.6\textwidth]{imagens/print_login_cliente.png}
\caption{Ecrã de autenticação do contribuinte (cidadão/comerciante).}
\label{fig:apendice_login_cliente}
\end{figure}

\begin{figure}[!htbp]
\centering
\includegraphics[width=0.85\textwidth]{imagens/print_registo.png}
\caption{Ecrã de registo de contribuinte.}
\label{fig:apendice_registo}
\end{figure}

\begin{figure}[!htbp]
\centering
\includegraphics[width=0.85\textwidth]{imagens/print_taxas_cidadao.png}
\caption{Painel do contribuinte com a lista de taxas disponíveis para pagamento.}
\label{fig:apendice_taxas_cidadao}
\end{figure}

\begin{figure}[!htbp]
\centering
\includegraphics[width=0.75\textwidth]{imagens/print_pagamento_emola.png}
\caption{Simulação de pagamento: seleção do método de pagamento via e-Mola.}
\label{fig:apendice_pagamento_emola}
\end{figure}

\begin{figure}[!htbp]
\centering
\includegraphics[width=0.75\textwidth]{imagens/print_erro_rede.png}
\caption{Validação de segurança: bloqueio de pagamento via e-Mola para um número Vodacom.}
\label{fig:apendice_erro_rede}
\end{figure}

\begin{figure}[!htbp]
\centering
\includegraphics[width=0.65\textwidth]{imagens/print_recibo.png}
\caption{Emissão do recibo digital após confirmação do pagamento.}
\label{fig:apendice_recibo}
\end{figure}

\section*{Portal do Administrador}
\begin{figure}[!htbp]
\centering
\includegraphics[width=0.6\textwidth]{imagens/print_login_admin.png}
\caption{Ecrã de autenticação do Administrador Municipal.}
\label{fig:apendice_login_admin}
\end{figure}

\begin{figure}[!htbp]
\centering
\includegraphics[width=0.9\textwidth]{imagens/print_admin_taxas.png}
\caption{Interface de gestão de taxas e categorias pelo administrador (UC03).}
\label{fig:apendice_admin_taxas}
\end{figure}

% IMAGEM ADICIONADA AQUI: CRIAR GESTOR
\begin{figure}[!htbp]
\centering
\includegraphics[width=0.9\textwidth]{imagens/print_admin_gestor.png}
\caption{Interface de criação de Gestor de Balcão pelo administrador.}
\label{fig:apendice_admin_gestor}
\end{figure}

% IMAGENS ADICIONADAS AQUI: RELATÓRIOS (DUAS IMAGENS)
\begin{figure}[!htbp]
\centering
\includegraphics[width=0.9\textwidth]{imagens/print_admin_relatorios.png}
\caption{Dashboard de arrecadação com filtros por data (UC04).}
\label{fig:apendice_admin_relatorios}
\end{figure}

\begin{figure}[!htbp]
\centering
\includegraphics[width=0.9\textwidth]{imagens/print_admin_relatorios_2.png}
\caption{Visualização detalhada de relatórios e exportação em CSV/PDF.}
\label{fig:apendice_admin_relatorios_2}
\end{figure}

\begin{figure}[!htbp]
\centering
\includegraphics[width=0.9\textwidth]{imagens/print_admin_auditoria.png}
\caption{Registo de auditoria do sistema com o histórico de ações e IP (UC07).}
\label{fig:apendice_admin_auditoria}
\end{figure}

\section*{Portal do Gestor de Balcão}
\begin{figure}[!htbp]
\centering
\includegraphics[width=0.6\textwidth]{imagens/print_login_gestor.png}
\caption{Ecrã de autenticação exclusivo do Gestor de Balcão.}
\label{fig:apendice_login_gestor}
\end{figure}

\begin{figure}[!htbp]
\centering
\includegraphics[width=0.9\textwidth]{imagens/print_gestor_balcao.png}
\caption{Janela do gestor de balcão para confirmação de pagamentos em numerário e POS (UC08).}
\label{fig:apendice_gestor_balcao}
\end{figure}

\section*{Evidências Técnicas (Backend)}
\begin{figure}[!htbp]
\centering
\includegraphics[width=0.9\textwidth]{imagens/print_backend_logs.png}
\caption{Terminal do backend a demonstrar o envio de SMS e e-mail.}
\label{fig:apendice_backend_logs}
\end{figure}

\end{document}